\section{Subquadratic integrability from the profile decomposition}
\label{subquad:section}

The strict finite-profile criterion gives more than absolute continuity:
the original joint distance law has a density in some \(L^p\) with
\(p>1\). This conclusion uses a variable deletion threshold. It does not
weaken the criterion on the dimensions. We first prove the elementary
norm estimate and then keep track of the discarded pin remainder.

\subsection{A parameterized decomposition estimate}
For a measurable function on a measure space \((X,m)\), use the weak
quasinorm
\[
 \|F\|_{p,\infty}=\sup_{t>0}t\,m\{|F|>t\}^{1/p}.
\]
\begin{lemma}\label{subquad:interpolation}
Suppose \(\|F\|_1\le1\), \(A>0\), \(\gamma>0\), \(K>0\), and,
for every \(\Lambda\ge1\), there is a decomposition
\[
 F=G_\Lambda+B_\Lambda,\qquad
 \|B_\Lambda\|_1\le\Lambda^{-\gamma},\qquad
 \|G_\Lambda\|_2^2\le A\Lambda^K.
\]
Put \(\theta=\gamma/(K+\gamma)\) and
\(p_0=1+\theta=(K+2\gamma)/(K+\gamma)\). Then
\begin{equation}\label{subquad:weak}
 \sup_{t>0}t^{p_0}m\{|F|>t\}\le6A^\theta,
 \qquad
 \|F\|_{p_0,\infty}\le6^{1/p_0}A^{\gamma/(K+2\gamma)}.
\end{equation}
If \(F\) is bounded, it suffices to have these decompositions for
\[
 1\le\Lambda\le\max\{1,(\|F\|_\infty/A)^{1/(K+\gamma)}\}.
\]
Fixed multiplicative constants in the assumptions change only the
conclusion's multiplicative constant.
\end{lemma}
\begin{proof}
For \(t\le A\), Markov gives
\(t^{p_0}m\{|F|>t\}\le t^{p_0-1}\le A^\theta\).
For \(t>A\), choose \(\Lambda=(t/A)^{1/(K+\gamma)}\). The triangle
inequality and Chebyshev give
\[
 m\{|F|>t\}\le4A\Lambda^Kt^{-2}+2\Lambda^{-\gamma}t^{-1}
 =6A^\theta t^{-p_0}.
\]
Values \(t>\|F\|_\infty\) have zero superlevel measure, proving the
finite-range assertion. Rescaling \(F\), and changing the constant
multiplying \(A\), handles fixed constants in the hypotheses.
\end{proof}

We also record two direct consequences of the distribution bound. If
\(F\) is supported in a region of measure \(V<\infty\), then, for
\(1\le p<p_0\),
\begin{equation}\label{subquad:strong-from-weak}
 \|F\|_p\le\left(\frac{p_0}{p_0-p}\right)^{1/p}
 V^{1/p-1/p_0}\|F\|_{p_0,\infty}.
\end{equation}
For any measurable \(A\) of finite measure,
\begin{equation}\label{subquad:restricted-set}
 \int_A|F|\dd m\le\frac{p_0}{p_0-1}
 \|F\|_{p_0,\infty}\,m(A)^{1-1/p_0}.
\end{equation}
Both follow by integrating the smaller of the trivial set-measure bound
and the weak distribution bound, splitting the level integral where they
are equal. Thus no interpolation or Lorentz-space theorem is needed.

\subsection{Varying the deletion thresholds}
\begin{theorem}[A stronger conclusion under the strict profile criterion]
\label{subquad:profile-theorem}
Assume the separated-probability hypotheses of the finite-profile measure
argument, and suppose that the bounded-chain estimate
\eqref{uf:chain} holds with \(\mathcal C<s-1\). Then the joint distance
law induced by \((y,x)\mapsto(y,|x-y|)\) on \(\nu\times\mu\) has a
density in \(L^p(d\nu\,dt)\) for some \(p>1\). In particular, its
pinned densities belong to \(L^p(dt)\) for \(\nu\)-almost every pin.

More precisely, make any fixed parameter choices in the strict proof
leaving a negative good-shell exponent \(-\kappa<0\), a chain-length
bound \(K\ge1\), and discarded pin mass at most
\(CR^{-\epsilon_{\rm reg}}\), with \(\epsilon_{\rm reg}>0\).
Let \(q>1\) be the radial-projection exponent and put
\[
 \gamma=\frac{q-1}{q},\qquad p_0=\frac{K+2\gamma}{K+\gamma}.
\]
The conclusion holds for every
\begin{equation}\label{subquad:p-range}
 1<p<\min\{p_0,1+\epsilon_{\rm reg}/2\}.
\end{equation}
\end{theorem}
\begin{proof}
At each frequency fix the finite regularization, its components, and all
spatial chains before changing any deletion thresholds. Let \(\Omega_R\)
be the union of the retained pin components. With \(L(1),H_{i,j}(1)\)
the original thresholds in \eqref{fprof:thresholds}, use
\begin{equation}\label{subquad:thresholds}
 L(\Lambda)=L(1)\Lambda^{1/q},\qquad
 H_{i,j}(\Lambda)=H_{i,j}(1)\Lambda,\qquad \Lambda\ge1.
\end{equation}
The source-heavy mass estimate \eqref{fprof:sourcebad} changes by the
factor \(\Lambda^{(1-q)/q}=\Lambda^{-\gamma}\). In
\eqref{fprof:pinbad}, the ratio of source and pin thresholds changes by
\(\Lambda^{1/q-1}=\Lambda^{-\gamma}\). Thus both deletion contributions
acquire the same factor. The source-heavy directions still serve only
as an auxiliary estimate; no new source cutoff is introduced.

There is one pin-threshold factor in each Fourier inflation step.
Consequently the squared good-shell estimate changes by at most
\(\Lambda^K\). Write
\[
 F_R(y,t)=\frac{d(d_y)_*(P_R\mu)}{dt}
\]
for the full smooth annular density. On \(\Omega_R\), the construction
therefore supplies a decomposition
\begin{gather}\label{subquad:regular-decomposition}
 F_R\mathbf1_{\Omega_R}=G_{R,\Lambda}+B_{R,\Lambda},\qquad
 \|G_{R,\Lambda}\|_2^2\le CR^{-\kappa}\Lambda^K,\\
 \|B_{R,\Lambda}\|_1\le CR^{-c}\Lambda^{-\gamma}
       +\operatorname{RapDec}(R),\qquad c>0.\nonumber
\end{gather}
Here and below the norms are on the joint pin-distance space. The component
weights sum to at most one, so there is no component-count loss.

For precision, these bounds are uniform when
\(1\le\Lambda\le R^M\), for any fixed \(M\), after choosing sufficiently
high orders in the localization estimates. All thresholds select subsets
of the same fixed standard packet family. Its size and its global derivative
bounds are polynomial in \(R\), uniformly over the subset. The initial
reconstruction, inherited identities, and selected-label embedding are
unchanged. The extra products of thresholds amplify an error by at most
\(R^{MK}\). Arbitrarily high integration-by-parts decay absorbs this
and every other fixed polynomial factor. Thus the error order can be chosen
after \(M\); there is no requirement of uniformity for superpolynomial
thresholds. Any residual rapid term in the good estimate is likewise
absorbed in its displayed main bound.

The entire annular multiplier has a uniformly bounded first kernel norm.
Pushforward contracts total variation, uniformly in the pin. Its spatial
kernel has supremum at most \(C_TR^2\). The source cutoff is fixed, and
the radii under consideration lie in a fixed bounded interval. Integration
of that bounded spatial density over a circle proves
\begin{equation}\label{subquad:full-shell-bounds}
 \sup_y\|F_R(y,\cdot)\|_1\le C_T,
 \qquad \|F_R\|_\infty\le C_TR^2.
\end{equation}
In particular the joint first norm of \(F_R\mathbf1_{\Omega_R}\) is
uniformly bounded.

Choose once and for all
\[
 M>\frac{2+\kappa}{K+\gamma}+1.
\]
The finite parameter range in Lemma~\ref{subquad:interpolation}, with
\(A=CR^{-\kappa}\), is then contained in \([1,R^M]\) for large \(R\).
Choose the rapid-error first norm in
\eqref{subquad:regular-decomposition} at most \(R^{-M\gamma-1}\).
Throughout the needed parameter range it is bounded by
\(\Lambda^{-\gamma}\). Applying the lemma, with fixed constants
absorbed, gives
\begin{equation}\label{subquad:regular-weak}
 \|F_R\mathbf1_{\Omega_R}\|_{p_0,\infty}
 \le CR^{-\kappa\gamma/(K+2\gamma)}.
\end{equation}

The discarded pins have to be treated separately. The good function is
zero there at every threshold, so their contribution does not gain a
factor \(\Lambda^{-\gamma}\). Their measure and
\eqref{subquad:full-shell-bounds} instead give
\[
 \|F_R\mathbf1_{\Omega_R^c}\|_1\le CR^{-\epsilon_{\rm reg}},
 \qquad
 \|F_R\mathbf1_{\Omega_R^c}\|_\infty\le CR^2.
\]
For \(p>1\), bound \(p-1\) factors of the integrand by its supremum
and integrate the remaining factor. This yields
\begin{equation}\label{subquad:remainder}
 \|F_R\mathbf1_{\Omega_R^c}\|_p
 \le C_pR^{[2(p-1)-\epsilon_{\rm reg}]/p}.
\end{equation}
For \(p\) in \eqref{subquad:p-range}, use
\eqref{subquad:strong-from-weak} on the common bounded radial support in
\eqref{subquad:regular-weak}, and add \eqref{subquad:remainder}. We obtain
\begin{gather}\label{subquad:shell-decay}
 \|F_R\|_p\le C_pR^{-c_p},\\
 c_p=\min\left\{\frac{\kappa\gamma}{K+2\gamma},
       \frac{\epsilon_{\rm reg}-2(p-1)}p\right\}>0.\nonumber
\end{gather}
The annular frequencies form a geometric sequence, so the shell norms
are summable. The series, including the smooth low-frequency density,
converges in \(L^p\), hence in \(L^1\) on the bounded joint support.

To identify the limit, test against a smooth function \(\Phi(y,t)\).
The spatial test
\[
 x\longmapsto\chi(x)\int\Phi(y,|x-y|)\dd\nu(y)
\]
is smooth and compactly supported because the source cutoff is separated
from the pin support. Distributional reconstruction of the annular
multipliers and first-norm convergence show that the limit is the joint
distance law. It is therefore a nonnegative density. Disintegration and
Fubini give the asserted almost-everywhere pinned conclusion.
\end{proof}

The exponent in this theorem depends on the regularization choices as well
as on \(q\) and \(K\). In particular, the weak exponent \(p_0\) proved
for the retained pin region is not automatically the exponent of the full
law: the discarded region imposes the second restriction in
\eqref{subquad:p-range}.

\subsection{The frequency threshold is unchanged}
Before using strictness, the leading squared good-shell factor is
\[
 A_R=R^{\mathcal C-(s-1)+\eta},\qquad \eta>0.
\]
Lemma~\ref{subquad:interpolation} raises this factor to a positive power.
The sign of its exponent is unchanged. Thus the present interpolation
closes under the same condition \(\mathcal C<s-1\).

The abstract estimate is sharp in its dependence on \(A\). For \(A\ge1\),
take \(F_A=A\mathbf1_{[0,1/A]}\). Its first norm is one and its squared
second norm is \(A\). Choosing the entire function as the good part and
zero as the bad part satisfies every hypothesis of the decomposition lemma,
whereas
\[
 \|F_A\|_{p_0,\infty}^{p_0}=A^{p_0-1}=A^{\gamma/(K+\gamma)}.
\]
This obstructs a better power from those hypotheses alone. It does not
exclude a different subquadratic estimate using additional geometry.

For comparison, a single-threshold decomposition with
\(\|B_H\|_1\le J/H\) and \(\|G_H\|_2^2\le HE\), \(H\ge J\),
gives \(p_0=3/2\) by putting \(\Lambda=H/J\). Its decisive factor is
\(A=JE\), so even this favorable case retains the original frequency sign.

\subsection{A limitation of higher pin-tube moments}
There is a similarly precise limitation to using higher moments estimated
only by the maximum tube mass. Consider an ideal regular edge from depth
\(n\) to depth \(m<n\), and put
\[
 \ell=n-m,\quad\rho=2^{-\ell},\quad
 c=g(m)-\min_{[m,n]}g,\quad
 d_* =\min\{\ell,\ell+g(m)-g(n)\}.
\]
Let \(\lambda\) be the normalized parent pin probability. A tube can be
covered by \(O(\rho^{-1})\) child-scale cubes. Regularity, together with
the trivial mass bound, gives
\[
 \frac{\max_T\lambda(T)}\rho\lesssim2^{d_*}.
\]
The pair bound gives \(J_2=\sum_T\lambda(T)^2\lesssim2^c\).
For \(r\ge2\), writing \(M=\max_T\lambda(T)\), one has
\[
 \sum_T\lambda(T)^r\le M^{r-2}J_2.
\]
On a retained direction the source mass is at most \(CL\rho\).
The higher-moment replacement in the heavy-tube calculation is therefore
\[
 \mathrm{BadMass}\lesssim
 L\rho(H\rho)^{1-r}\sum_T\lambda(T)^r
 \lesssim LH^{-(r-1)}2^{c+(r-2)d_*}.
\]
To get a small bad mass, this uses the base threshold exponent
\[
 c_r=\frac{c+(r-2)d_*}{r-1}.
\]
Both entries defining \(d_*\) are at least \(c\). The first follows
from the Lipschitz bound. If \(z\) is a minimum point on the edge, then
\[
 \ell+g(m)-g(n)-c
 =\ell+g(z)-g(n)\ge\ell-(n-z)\ge0.
\]
Consequently \(c_r\ge c\). The improved tail power is accompanied by a
base inflation exponent no smaller than before. The finite regularization
and geometric enlargement errors add the previously accounted small powers;
they do not reverse this comparison. This limits this particular use of
higher moments, not every possible higher-moment argument.
